%% ==========================================================================
%%  Approach 2 --- select a utilitarian-optimal allocation.   DEAD.
%%
%%  Presentation version: the idea, one counterexample, the verdict.  The
%%  tie-break post-mortems and the sampler-bias discussion are removed; they
%%  are in git history.
%% ==========================================================================
\section{Approach 2: Selecting a Utilitarian-Optimal Allocation}
\label{sec:approach2}

\paragraph{The idea.} Theorem~\ref{thm:obstruction} rules out building the
allocation one chore at a time. The natural response is to stop building it and
instead \emph{select} one globally, from a set small enough to reason about and
rich enough to contain a witness. Utilitarian optimality is the obvious
candidate.

\paragraph{Why it looks promising.} Two reasons, both real. First, by
Theorem~\ref{thm:hs-characterisation} an allocation minimising $\sum_i
\cost_i(\bundle{i})$ over \emph{all} allocations --- not merely over
reassignments of its own bundles --- is automatically envy-freeable, so half of
what is needed comes for free. Second, it carries a strong exchange property: if
$g \in \bundle{i}$ is \emph{active} for $i$, meaning $\cost_i(\bundle{i}) -
\cost_i(\bundle{i} \setminus \set{g}) = 1$, then $\cost_j(\bundle{j} \cup
\set{g}) = \cost_j(\bundle{j}) + 1$ for every $j$ --- nobody can absorb it for
free, or the allocation would not have been optimal. With
Theorem~\ref{thm:cyclebound}, a path of weight at least $2$ from $u$ to $v$
forces $\cost_v(\bundle{u}) \ge \cost_v(\bundle{v}) + 2$, so $\bundle{u}$ carries
two units live for $v$, and one would like to transfer one along the path and
watch a potential drop.

\paragraph{Where it fails.}

\begin{proposition}
\label{prop:msw-false}
There is an instance with $n = 3$ and $m = 4$ whose \emph{unique}
utilitarian-optimal allocation requires a subsidy of $2$, although another
allocation --- of strictly greater total cost --- admits a subsidy of $0$.
\end{proposition}

\begin{proof}
Take $\items = \set{g_1,g_2,g_3,g_4}$ with the costs below.

\begin{center}\small
\begin{tabular}{@{}lccc@{\hspace{2.5em}}lccc@{}}
\toprule
$S$ & $\cost_1$ & $\cost_2$ & $\cost_3$ &
$S$ & $\cost_1$ & $\cost_2$ & $\cost_3$ \\
\midrule
$\emptyset$      & 0 & 0 & 0 & $g_3g_4$       & 2 & 1 & 1 \\
$g_1$            & 1 & 1 & 1 & $g_1g_2g_3$    & 2 & 3 & 2 \\
$g_2$            & 1 & 1 & 1 & $g_1g_2g_4$    & 3 & 2 & 2 \\
$g_3$            & 1 & 1 & 0 & $g_1g_3g_4$    & 2 & 2 & 2 \\
$g_4$            & 1 & 1 & 1 & $g_2g_3g_4$    & 2 & 2 & 2 \\
$g_1g_2$         & 2 & 2 & 2 & $g_1g_2g_3g_4$ & 3 & 3 & 3 \\
$g_1g_3$         & 1 & 2 & 1 & & & & \\
$g_1g_4$         & 2 & 2 & 2 & & & & \\
$g_2g_3$         & 1 & 2 & 1 & & & & \\
$g_2g_4$         & 2 & 2 & 2 & & & & \\
\bottomrule
\end{tabular}
\end{center}

Each column is $0$ at $\emptyset$ with every marginal in $\set{0,1}$, so all
three are negative dichotomous in cost form.

Agent~1 pays $1$ for every singleton, so $\cost_1(\bundle{1}) = 0$ only if
$\bundle{1} = \emptyset$; agent~3 pays $0$ for $\set{g_3}$ but $1$ for every
superset. Any allocation of total cost below $2$ would need two agents at cost
$0$, forcing $\bundle{2} \supseteq \set{g_1,g_2,g_4}$, which costs agent~2 at
least $2$. So the minimum total cost is $2$, attained only by
\[
  \alloc^{\star} = \big(\emptyset,\ \set{g_1,g_2,g_4},\ \set{g_3}\big),
  \qquad 0 + 2 + 0 = 2,
\]
whose envy graph has $\edgew{\alloc^\star}(2,1) = 2$, giving $\optsubsidy =
(0,2,0)$. But $\allocB = (\set{g_1,g_3},\ \set{g_2},\ \set{g_4})$, of total cost
$3$, has envy matrix
\[
  \big(\edgew{\allocB}(i,j)\big) =
  \begin{pmatrix} 0 & 0 & 0 \\ -1 & 0 & 0 \\ 0 & 0 & 0 \end{pmatrix},
\]
every entry at most $0$, hence $\optsubsidy = (0,0,0)$.
\end{proof}

Two features make this sharper than a bare counterexample: the optimum is
\emph{unique}, so there is no tie to break, and the gap is the largest possible
in one step --- subsidy $2$ against $0$, for one extra unit of total burden.

\begin{corollary}[No refinement of the utilitarian-optimal set can work]
\label{cor:no-refinement}
Let $\mathcal{R}$ be any rule selecting a nonempty subset of the
utilitarian-optimal allocations. Then $\mathcal{R}$ fails on the instance of
Proposition~\ref{prop:msw-false}: the optimal set is the singleton
$\set{\alloc^\star}$, so $\mathcal{R}$ returns it, and
$\pathw{\alloc^\star}(2) = 2$ while a subsidy-free allocation exists outside.
\end{corollary}

\begin{remark}[Verdict]
\label{rem:msw-deadend}
Conjecture~\ref{conj:main}, if true, is \emph{not} witnessed by welfare-maximising
allocations, and looking for a tie-breaking rule inside the utilitarian-optimal
set is unsound as posed. The exchange property quoted above remains true and
usable; what fails is the assumption that a witness may be sought among
cost-minimising allocations at all. \textbf{Any search procedure must be free to
give up total cost.} The slip underneath is a misreading of
Theorem~\ref{thm:hs-characterisation}: it says an envy-freeable allocation
maximises welfare across reassignments of \emph{its own bundles}, not that it is
globally welfare-maximising --- and $\allocB$ above, envy-freeable at cost $3$
against the optimum $2$, is a direct witness that the two notions differ.
\end{remark}

\begin{remark}[The refinement that looked safe]
\label{rem:double-refinement}
Refining the optimal set lexicographically by cardinality balance and then by
own-cost --- suggested independently by Theorem~\ref{thm:sizeshift},
by~\cite{BDNSV20}'s balancedness guarantee, and by Remark~\ref{rem:balance} ---
survives randomised search over roughly $1050$ instances and is nevertheless
false by Corollary~\ref{cor:no-refinement}, both refinements being no-ops on a
singleton. The failure rate is about one instance in $2.6 \times 10^4$, so the
sweep was two orders of magnitude short of seeing it (\Cref{sec:samplerbias}).
\end{remark}

\subsection{Selection rules that fail}
\label{sec:tiebreaks}

Three tie-breaking rules inside the utilitarian-optimal set were tested before
Proposition~\ref{prop:msw-false} was found --- minimising total perceived load,
leximin on the perceived-load vector, and local search by single transfers --- and
all three fail. The informative one is the last: single transfers are too coarse,
because the utilitarian-optimal allocations of a dichotomous instance have
natural moves that are exchange \emph{paths and cycles}, as Yankee Swap and
matroid union are in the matroid-rank literature~\cite{BEF21}. Exchange-path local
search over the full allocation lattice --- not the optimal set, per
Remark~\ref{rem:msw-deadend} --- remains untried.
